% GENERATED FILE. Edit canonical lessons, then run npm run generate:books.
% canonical-source-hash: fnv1a64:cf3327bf20cf7a5e
% canonical-lessons: ES-C591-convenio, ES-C591-autonomo, ES-C591-funcionario, ES-C591-teletrabajo, ES-C591-sindicato, ES-R591-repaso-584-591

\chapter{Formas de trabajar}
\label{ch:es-formas-de-trabajar}
\hlchaptermodality{591}

\begin{chapteropening}
\textbf{By the end of this chapter:} \emph{I can talk about civil servants, the self-employed, working from home, the union and its agreement.}

Complete the last lesson of chapter 591: I can talk about civil servants, the self-employed, working from home, the union and its agreement.
\end{chapteropening}

\section[el convenio]{el convenio — the agreement; the collective agreement}

\label{lesson:ES-C591-convenio}



\begin{quote}
Before the new one: say the Spanish for the training, then the Spanish for the profile.
\end{quote}

\subsection*{You'll want to know: el convenio}
\textbf{el convenio} — \textquotedblleft{}the agreement; the collective agreement\textquotedblright{}.

\begin{quote}
El convenio sube los sueldos un dos por ciento.
\end{quote}

\begin{quote}
\emph{The collective agreement raises wages by two per cent.}
\end{quote}

\emph{El convenio} sets pay and hours for everyone in one kind of work.

\begin{cousinweb}
\emph{convenio} comes from \emph{convenir}, inherited from Latin \emph{convenire} 'to come together, to agree' (EtymDB), built on \emph{venire} 'to come'.
\end{cousinweb}

\begin{grammarlens}[title={What you've built}]
One more word for telling, explaining and describing at B1.
\end{grammarlens}

\subsection*{Your turn}
\begin{itemize}
\raggedright
  \item \emph{Say it:} \emph{el convenio}
  \item \emph{Say it:} \emph{el convenio}, once more
  \item \emph{Say it:} a sentence of your own with \emph{el convenio}
  \item \emph{Recall:} say the Spanish for the training, then the Spanish for the profile, then say \emph{el convenio} again
\end{itemize}

\begin{tcolorbox}[breakable,colback=teal!4,colframe=teal!35!black,title={Before you move on}]
What does el convenio mean? (\textquotedblleft{}The agreement; the collective agreement\textquotedblright{}.) Say it once more.
\end{tcolorbox}
\section[el autónomo]{el autónomo — the self-employed worker; autonomous}

\label{lesson:ES-C591-autonomo}



\begin{quote}
Before the new one: say the Spanish for the profile, then the Spanish for the agreement.
\end{quote}

\subsection*{You'll want to know: el autónomo}
\textbf{el autónomo} — \textquotedblleft{}the self-employed worker; autonomous\textquotedblright{}.

\begin{quote}
Mi hermana trabaja como autónoma.
\end{quote}

\begin{quote}
\emph{My sister is self-employed.}
\end{quote}

In Spain \emph{ser autónomo} is to be self-employed; much of Latin America says \emph{trabajador independiente}.

\begin{cousinweb}
\emph{autónomo} is a learned borrowing from Greek \emph{autonomos} 'having its own laws', built from \emph{autos} 'self' and \emph{nomos} 'custom, law' (EtymDB).
\end{cousinweb}

\begin{grammarlens}[title={What you've built}]
One more word for telling, explaining and describing at B1.
\end{grammarlens}

\subsection*{Your turn}
\begin{itemize}
\raggedright
  \item \emph{Say it:} \emph{el autónomo}
  \item \emph{Say it:} \emph{el autónomo}, once more
  \item \emph{Say it:} a sentence of your own with \emph{el autónomo}
  \item \emph{Recall:} say the Spanish for the profile, then the Spanish for the agreement, then say \emph{el autónomo} again
\end{itemize}

\begin{tcolorbox}[breakable,colback=teal!4,colframe=teal!35!black,title={Before you move on}]
What does el autónomo mean? (\textquotedblleft{}The self-employed worker; autonomous\textquotedblright{}.) Say it once more.
\end{tcolorbox}
\section[el funcionario]{el funcionario — the civil servant, the public employee}

\label{lesson:ES-C591-funcionario}



\begin{quote}
Before the new one: say the Spanish for the agreement, then the Spanish for the self-employed worker.
\end{quote}

\subsection*{You'll want to know: el funcionario}
\textbf{el funcionario} — \textquotedblleft{}the civil servant, the public employee\textquotedblright{}.

\begin{quote}
Mi madre es funcionaria del ayuntamiento.
\end{quote}

\begin{quote}
\emph{My mother is a civil servant at the town council.}
\end{quote}

A woman is \emph{la funcionaria}. In Spain many people sit public exams to become one.

\begin{cousinweb}
\emph{funcionario} is built on \emph{función}, a learned borrowing from Latin \emph{functio} 'a performance, a duty' (EtymDB), from \emph{fungi} 'to carry out'.
\end{cousinweb}

\begin{grammarlens}[title={What you've built}]
One more word for telling, explaining and describing at B1.
\end{grammarlens}

\subsection*{Your turn}
\begin{itemize}
\raggedright
  \item \emph{Say it:} \emph{el funcionario}
  \item \emph{Say it:} \emph{el funcionario}, once more
  \item \emph{Say it:} a sentence of your own with \emph{el funcionario}
  \item \emph{Recall:} say the Spanish for the agreement, then the Spanish for the self-employed worker, then say \emph{el funcionario} again
\end{itemize}

\begin{tcolorbox}[breakable,colback=teal!4,colframe=teal!35!black,title={Before you move on}]
What does el funcionario mean? (\textquotedblleft{}The civil servant, the public employee\textquotedblright{}.) Say it once more.
\end{tcolorbox}
\section[el teletrabajo]{el teletrabajo — remote work, working from home}

\label{lesson:ES-C591-teletrabajo}



\begin{quote}
Before the new one: say the Spanish for the self-employed worker, then the Spanish for the civil servant.
\end{quote}

\subsection*{You'll want to know: el teletrabajo}
\textbf{el teletrabajo} — \textquotedblleft{}remote work, working from home\textquotedblright{}.

\begin{quote}
Hago teletrabajo dos días a la semana.
\end{quote}

\begin{quote}
\emph{I work from home two days a week.}
\end{quote}

The verb is \emph{teletrabajar}.

\begin{cousinweb}
\emph{teletrabajo} is built on Greek \emph{tele} 'far' and \emph{trabajo}. \emph{Trabajo} is usually traced to Vulgar Latin \emph{tripālium}, a device for torture (EtymDB), so the word for work first meant suffering.
\end{cousinweb}

\begin{grammarlens}[title={What you've built}]
One more word for telling, explaining and describing at B1.
\end{grammarlens}

\subsection*{Your turn}
\begin{itemize}
\raggedright
  \item \emph{Say it:} \emph{el teletrabajo}
  \item \emph{Say it:} \emph{el teletrabajo}, once more
  \item \emph{Say it:} a sentence of your own with \emph{el teletrabajo}
  \item \emph{Recall:} say the Spanish for the self-employed worker, then the Spanish for the civil servant, then say \emph{el teletrabajo} again
\end{itemize}

\begin{tcolorbox}[breakable,colback=teal!4,colframe=teal!35!black,title={Before you move on}]
What does el teletrabajo mean? (\textquotedblleft{}Remote work, working from home\textquotedblright{}.) Say it once more.
\end{tcolorbox}
\section[el sindicato]{el sindicato — the trade union}

\label{lesson:ES-C591-sindicato}



\begin{quote}
Before the new one: say the Spanish for the civil servant, then the Spanish for remote work.
\end{quote}

\subsection*{You'll want to know: el sindicato}
\textbf{el sindicato} — \textquotedblleft{}the trade union\textquotedblright{}.

\begin{quote}
El sindicato pide salarios más altos.
\end{quote}

\begin{quote}
\emph{The union is asking for higher wages.}
\end{quote}

A trade unionist is \emph{un sindicalista}; a member is \emph{un afiliado}.

\begin{cousinweb}
\emph{sindicato} is built on \emph{síndico} 'representative', a learned borrowing from Latin \emph{syndicus} and Greek \emph{syndikos} 'advocate' (EtymDB).
\end{cousinweb}

\begin{grammarlens}[title={What you've built}]
One more word for telling, explaining and describing at B1.
\end{grammarlens}

\subsection*{Your turn}
\begin{itemize}
\raggedright
  \item \emph{Say it:} \emph{el sindicato}
  \item \emph{Say it:} \emph{el sindicato}, once more
  \item \emph{Say it:} a sentence of your own with \emph{el sindicato}
  \item \emph{Recall:} say the Spanish for the civil servant, then the Spanish for remote work, then say \emph{el sindicato} again
\end{itemize}

\begin{tcolorbox}[breakable,colback=teal!4,colframe=teal!35!black,title={Before you move on}]
What does el sindicato mean? (\textquotedblleft{}The trade union\textquotedblright{}.) Say it once more.
\end{tcolorbox}
\section[(dialogue)]{Repaso de los capítulos 584–591}

\label{lesson:ES-R591-repaso-584-591}



\begin{quote}
Say the Spanish for the advantage and the disadvantage.
\end{quote}

\subsection*{Your turn}
Say these aloud:

\begin{itemize}
\raggedright
  \item the advantage, the disadvantage, the drawback, positive, negative
  \item to propose, the proposal, the suggestion, to warn, the warning
  \item to justify, the justification, the excuse, the pretext, to show
  \item to check, evident, obvious, the certainty, apparently
  \item therefore, I mean, whereas, above all, in general
  \item in fact, actually, certainly, not at all, deep down
  \item the vacancy, to select, the employee, the training, the profile
  \item the agreement, the self-employed worker, the civil servant, remote work, the trade union
\end{itemize}

\begin{tcolorbox}[breakable,colback=teal!4,colframe=teal!35!black,title={Before you move on}]
The advantage? (\textbf{la ventaja}.) The disadvantage? (\textbf{la desventaja}.) The drawback? (\textbf{el inconveniente}.) Positive? (\textbf{positivo}.) Negative? (\textbf{negativo}.) To propose? (\textbf{proponer}.) The proposal? (\textbf{la propuesta}.) The suggestion? (\textbf{la sugerencia}.) To warn? (\textbf{advertir}.) The warning? (\textbf{la advertencia}.) To justify? (\textbf{justificar}.) The justification? (\textbf{la justificación}.) The excuse? (\textbf{la excusa}.) The pretext? (\textbf{el pretexto}.) To show? (\textbf{demostrar}.) To check? (\textbf{comprobar}.) Evident? (\textbf{evidente}.) Obvious? (\textbf{obvio}.) The certainty? (\textbf{la certeza}.) Apparently? (\textbf{aparentemente}.) Therefore? (\textbf{por tanto}.) I mean? (\textbf{o sea}.) Whereas? (\textbf{en cambio}.) Above all? (\textbf{sobre todo}.) In general? (\textbf{en general}.) In fact? (\textbf{de hecho}.) Actually? (\textbf{en realidad}.) Certainly? (\textbf{desde luego}.) Not at all? (\textbf{en absoluto}.) Deep down? (\textbf{en el fondo}.) The vacancy? (\textbf{la vacante}.) To select? (\textbf{seleccionar}.) The employee? (\textbf{el empleado}.) The training? (\textbf{la formación}.) The profile? (\textbf{el perfil}.) The agreement? (\textbf{el convenio}.) The self-employed worker? (\textbf{el autónomo}.) The civil servant? (\textbf{el funcionario}.) Remote work? (\textbf{el teletrabajo}.) The trade union? (\textbf{el sindicato}.)
\end{tcolorbox}
